%% Para compilar o programa, use pdflatex
%% Opcoes:
%%
%% Cursos:
%% cic: Bacharelado em Ciência da Computação
%% tsi: Técnologo em Sistemas para Internet
\documentclass[tsi]{ifbclass/ifbclass}

\address{BRASÍLIA}

\title{Ferramenta de Desenvolvimentos de Software na Web para aprendizagem de algoritmos}

\date{2019}

\author{Neander Wendel Nobre Teixeira}
\adviser{Prof.º Alisson Wilker Andrade Silva}

\membroum[a]{Dr.ª Primeira Membro da Banca}
\membrodois{Dr. Segundo Membro da Banca}
\membrotres[a]{Dr.ª Terceira Membro da Banca}
%\membroquatro[a]{Dr.ª Quarta Membro da Banca}

%\coadviser{Nome completo do co-orientador (se existir)} %TODO: nao suportado

% Macros (Se necessario, defina suas proprias aqui)
\def\x{\checkmark}
%\let\lstlistoflistings\origlstoflistings
  \begin{document}

\frontmatter

\frontpage

\presentationpage

% Quando a biblioteca preparar a ficha catalografica, insira-a aqui
% Antes da defesa do trabalho, crie uma ficha falsa
\begin{fichacatalografica}
  \FakeFichaCatalografica
%     \includepdf{fig_ficha_catalografica.pdf} % crie o arquivo come esse nome e descomente essa linha
\end{fichacatalografica}

\banca

%Comente/Remova todo o conteudo entre \begin{dedicatory} e \end{dedicatory} para remover.
\begin{dedicatory} %OPCIONAL
Dedico este trabalho à minha família e a todos que me apoiaram.
\end{dedicatory}
  
\acknowledgements %OPCIONAL Comente a linha abaixo para remover.
Agradeço ao meu orientador Prof. Dr. Nome do Orientador, pela sabedoria com que me guiou nesta trajetória.

Aos meus colegas de sala.

A Secretaria do Curso, pela cooperação.

Gostaria de deixar registrado também, o meu reconhecimento à minha família, pois acredito que sem o apoio deles seria muito difícil vencer esse desafio. 

Enfim, a todos os que por algum motivo contribuíram para a realização desta pesquisa.


%Novamente, comente todo conteudo entre {epigraph} para remover
\begin{epigraph}[]{Nome do autor} %OPCIONAL
Elemento opcional.

Espaço destinado à epígrafe (elemento opcional). Nesta folha, o autor usa uma citação, seguida de indicação de autoria e ano, relacionada com a matéria tratada no corpo do trabalho.
\end{epigraph}

\resumo
% Escreva seu resumo no arquivo resumo.tex
{\parindent0pt
  SOBRENOME, Prenome do Autor do Trabalho. Título do trabalho: subtítulo (se houver).  2018. 65 f. 
Trabalho de Conclusão de Curso (Graduação) – Tecnólogo em Sistemas para Internet. 
Instituto Federal de Brasília – Campus Brasília. Brasília/DF, 2018.
\vspace{1cm}

Elemento obrigatório, constituído de uma sequência de frases concisas e objetivas,
fornecendo uma visão rápida e clara do conteúdo do estudo. O texto deverá conter no
máximo 500 palavras e ser antecedido pela referência do estudo, com exceção do resumo
inserido no próprio documento. Também, não deve conter citações. O resumo deve ser redigido
em parágrafo único, espaçamento simples e seguido das palavras representativas do conteúdo
do estudo, isto é, palavras-chave, em número de três a cinco, separadas entre si por ponto e
finalizadas também por ponto. Usar o verbo na terceira pessoa do singular, com linguagem
impessoal (pronome SE), bem como fazer uso, preferencialmente, da voz ativa.

\begin{keywords}
Primeira palavra. Segunda palavra. Terceira palavra. Quarta palavra. Quinta-palavra.
\end{keywords}

}
  
\abstract
% Escreva seu abstract em ingles no arquivo abstract.tex
{\parindent0pt
  SOBRENOME, Prenome do Autor do Trabalho. Título do trabalho: subtítulo (se houver).  2018. 65 f. 
Trabalho de Conclusão de Curso (Graduação) – Tecnólogo em Sistemas para Internet. 
Instituto Federal de Brasília – Campus Brasília. Brasília/DF, 2018.
\vspace{1cm}

Elemento obrigatório, constituído de uma sequência de frases concisas e objetivas,
fornecendo uma visão rápida e clara do conteúdo do estudo. O texto deverá conter no
máximo 500 palavras e ser antecedido pela referência do estudo, com exceção do resumo
inserido no próprio documento. Também, não deve conter citações. O resumo deve ser redigido
em parágrafo único, espaçamento simples e seguido das palavras representativas do conteúdo
do estudo, isto é, palavras-chave, em número de três a cinco, separadas entre si por ponto e
finalizadas também por ponto. Usar o verbo na terceira pessoa do singular, com linguagem
impessoal (pronome SE), bem como fazer uso, preferencialmente, da voz ativa.

\begin{keywords}
Keyword. Second keyword. Third keyword. Keyword.
\end{keywords}

}

%ATENCAO.
%Se alguma dessas listas estiver vazia no seu trabalho, comente a linha com um '%'
% Lista de figuras
\listoffigures

% Lista de algortimos
\lstlistoflistings

% Lista de tabelas
\listoftables

% Lista de acronimos
% Acronyms manual: http://linorg.usp.br/CTAN/macros/latex/contrib/acronym/acronym.pdf
\listofacronyms
\begin{acronym}[ACRONYM] 
% Change the word ACRONYM above to change the acronym column width.
% The column width is equals to the width of the word that you put.
% Read the manual about acronym package for more examples:
%   http://linorg.usp.br/CTAN/macros/latex/contrib/acronym/acronym.pdf
\acro{afm}[AFM]{Alphabet Frequency Matrix}
\acro{api}[API]{Application Programming Interface}
\acro{arima}[ARIMA]{Auto-Regressive Integrated Moving Average}
\acro{brn}[BRN]{Bug Report Network}
\acro{bts}[BTS]{Bug Triage System}
\acro{cas}[CAS]{Context-Aware Systems}
\acro{ccb}[CCB]{Change Control Board}
\acro{cr}[CR]{Change Request}
\acro{cvs}[CVS]{Concurrent Version System}
\acro{es}[ES]{Expert System}
\acro{floss}[FLOSS]{Free/Libre Open Source Software}
\acro{glr}[GLR]{Generalized Linear Regression}
\acro{gqm}[GQM]{Goal Question Metric}
\acro{html}[HTML]{HyperText Markup Language}
\acro{ir}[IR]{Information Retrieval}
\acro{irt}[IRT]{Recôncavo Institute of Technology}
\acro{jdt}[JDT]{Jazz Duplicate Finder}
\acro{lda}[LDA]{Latent Dirichlet Allocation}
\acro{loc}[LOC]{Lines of Code}
\acro{lsi}[LSI]{Latent Semantic Indexing}
\acro{ms}[MS]{Mapping Study}
\acro{msr}[MSR]{Mining Software Repositories}
\acro{nlp}[NLP]{Natural Language Processing}
\acro{promise}[PROMISE]{Predictive Models in Software Engineering}
\acro{rbes}[RBES]{Rule-Based Expert System}
\acro{rhel}[RHEL]{RedHat Enterprise Linux}
\acro{saas}[SaaS]{Software as a Service}
\acro{scm}[SCM]{Software Configuration Management}
\acro{serpro}[SERPRO]{Brazilian Federal Organization for Data Processing}
\acro{slr}[SLR]{Stepwise Linear Regression}
\acro{slr}[SLR]{Systematic Literature Review}
\acro{svd}[SVD]{Singular Value Decomposition}
\acro{svm}[SVM]{Support Vector Machine}
\acro{svn}[SVN]{Subversion}
\acro{tfidf}[TF-IDF]{Term Frequency-Inverse Document Frequency}
\acro{vsm}[VSM]{Vector Space Model}
\acro{xp}[XP]{Extreming Programming}
\end{acronym}

% Sumario
\tableofcontents

\mainmatter

%Para cada seção do seu trabalho, edite o arquivo .tex na pasta text/
\chapter{Introdução}
A área de Desenvolvimento de Sistemas é uma área relativamente nova, mas que atualmente impacta em toda e qualquer área do conhecimento. Desde pequenos sistemas de gerenciamento de tarefas a grandes redes sociais, os Sistemas de Informação estão em cada momento de nossas vidas e inclusive, na educação.

\section{Tema}
As ferramentas disponíveis na Web estão cada vez mais sofisticadas e completas, e dispensando o uso de papéis ou de programas instalados fisicamente em suas máquinas, além de possuírem armazenamento compartilhado em nuvem. A utilização dessas ferramentas trazem mais praticidade e conforto ao usuário, principalmente para estudantes por manter todos os seus trabalhos, apresentações, anotações salvas na nuvem e acessíveis por qualquer dispositivo.

\section{Problema}
Atualmente, quando se começa a aprender uma linguagem de programação, os instrutores geralmente começam com uma pseudo-linguagem para ensinar os básicos de algoritmos. Em geral, os ambientes dessas pseudo-linguagens são simples e fáceis de usar, com apenas alguns cliques.

Entretanto, quando passa-se para uma linguagem de programação real (como Java, C\#, Python, etc...) sempre há problemas na primeira vez de executar os códigos. Variáveis de ambiente, máquina mal ambientada ou problemas de hardwares são alguns dos fatores que podem atrapalhar na hora de estudar uma linguagem de programação robusta.

\subsection{Objetivo geral}
Visando facilitar o início da aprendizagem com uma nova linguagem de programação uma ferramenta Web pode ajudar nessa transição para evitar que o estudante desista ainda na ambientação, além de motivá-lo a continuar estudando algoritmos sem precisar ter uma máquina potente ou um ambiente instalado em qualquer lugar.

\subsection{Objetivos específicos}
O sistema deverá:
\begin{itemize}
	\item Permitir ao usuário escolher a linguagem de programação desejada para resolver os algoritmos.
	\item Comparar suas estatísticas com a de outros usuários, com o intuito de criar um pequeno ambiente de competição.
	\item Possibilitar o usuário a desenvolver seus algoritmos sem preocupar-se com a ambientação da linguagem.
	\item Disponibilizar uma forma de criar um Avatar para sentir-se mais próximo dentro da plataforma.
\end{itemize}


\section{Estrutura do TCC}
Este trabalho consistirá em uma documentação para a Ferramenta de Desenvolvimento de Software na Web.
Inicialmente haverá uma breve revisão de sistemas similares disponíveis no mercado e logo após se iniciará a documentação do sistema, com o Documento de Visão, Descrição resumida dos casos de uso, diagramas de caso de uso, descrição detalhada dos casos de uso, protótipos de interfaces do sistema, diagramas de classe, diagramas de sequência e casos de teste.
Alguns desses documentos estarão em outros arquivos fora deste, cada seção definirá como será descrita.

\subsection{Classificação da Pesquisa}
Esta pesquisa será descritiva onde procurará identificar os fatores que envolvem os estudantes e a forma como eles aprendem. Também tem como objetivo documentar e explicar como o sistema será desenvolvido.
Será desenvolvido um sistema Web, utilizando a biblioteca React JS para as interfaces de usuário e a plataforma ASP.NET Core 2.2 para desenvolvimento dos serviços disponibilizados.



\chapter{Entrega parcial}
Os outros capítulos ainda estão em desenvolvimento, esta é apenas uma entrega parcial.

\chapter{Conceitos gerais e revisão da literatura}
\label{chp:conceitos-gerais}
Conforme explorado no capítulo \ref{chp:introducao}, um ambiente de desenvolvimento na Web é essencial para o engajamento de estudantes na área de desenvolvimento de sistemas. Com isso, nossa ferramenta Web de desenvolvimento pretende disponibilizar algoritmos e um ambiente de desenvolvimento on-line, com a possibilidade de salvar os algoritmos na núvem e permitirá validar eles. Os usuários também poderão analisar onde estão acertando e errando para que possam se auto-avaliar.

A disponibilização de ambientes de desenvolvimento na Web não é novidade. Entretanto, cada ferramenta tem sua especialidade e foco em uma área de desenvolvimento. Neste capítulo serão apresentados sistemas que seguem essa temática e que inspiram essa ferramenta.

\section{CodePen}
\label{sec:codepen}
O CodePen é um ambiente de desenvolvimento colaborativo com foco no desenvolvimento de Front-end e de código aberto. O CodePen foi desenvolvido como uma comunidade de compartilhamento de demonstrações e exemplos. Ele permite que os usuários desenvolvam sites através do navegador. Ele é ideal para expor trabalhos interativos e funcionais \cite{fiala2016collaborative}.

\section{JSFiddle}
\label{sec:jsfiddle}
O JSFiddle também é um ambiente de desenvolvimento colaborativo. Tal qual o CodePen \ref{sec:codepen}, também permite o compartilhamento de demonstrações e exemplos, com uma comunidade ampla. Possui suporte para frameworks e/ou bibliotecas mais robustas como React JS e VueJS.

\section{Coding Ground}
\label{sec:codingground}
O Coding Ground é uma plataforma de desenvolvimento para sites de tutoriais. Ele disponibiliza um pequeno ambiente para que o usuário possa editar códigos e executá-los. Também possui uma vasta biblioteca de linguagens possíveis para utilização, inclusive com frameworks e bibliotecas de Front-End.

\section{Fabriki}
\label{sec:fabriki}
Uma plataforma educacional com foco em ensinar linguagem de programação. Criada para guiar o estudante do iniciante ao avançado. Possui suporte para desenvolvimento Web e algoritmos em Java. A plataforma também possibilita a validação dos algoritmos quando submetidos. É uma das inspirações para a Program.Acad.

\section{The Huxley}
\label{sec:thehuxley}
Uma plataforma competitiva, que possibilita o desenvolvimento direto pelo Browser, com validação de códigos e algoritmos desafiadores. Esta plataforma possui suporte para  até 7 linguagens de programação que o usuário pode escolher ao começar. Está disponível gratuitamente. Também é uma das inspirações para a Program.Acad.

%Neste capítulo deve ser proporcionado o estado da arte / referencial teórico
%sobre o tema a que se refere o estudo. Um bom pesquisador não deve repetir
%trabalhos já concluídos ou que já estão em andamento. Por isso esta sessão é
%onde o autor demonstra até onde vai a pesquisa atual no campo de estudos em
%questão e estabelece as bases sobre as quais desenvolverá o estudo proposto. A
%seguir são mostrados alguns exemplos de como deve-se inserir as figuras e
%tabelas. A Figura \ref{fig:exemplo} mostra um exemplo de como inserir uma
%figura no texto. A Tabela \ref{tb:exemplo} mostra o exemplo de como uma tabela
%deve ser inserida.  Voce pode referenciar capítulos e seções adicionando labels
%à elas. Por exemplo, descrevemos a introdução no Capítulo
%\ref{chp:introducao}.

\begin{figure}[htb]
    \centering
    \includegraphics[width=0.45\textwidth]{images/figura.png}
    \caption{Exemplo de como inserir Figura}
    \label{fig:exemplo}
\end{figure}
%Altere o numero junto ao 'width' para acertar o tamanho da figura. Coloque
%numeros maior que 0 e até 1

\begin{table}[htb]
\caption{Modelo de como as tabelas devem ser inseridas no texto}
\label{tb:exemplo}
\centering
\begin{tabular}{|l|c|r|r|} %left, center, right. Você pode mudar isso
\hline
Índice  & Coluna 1 & Coluna 2 & Coluna 3 \\
\hline
Linha 1 &          &          &          \\
Linha 2 &          &          &          \\
Linha 3 &          &          &          \\
\hline
\end{tabular}
\end{table}
%Para manter a sanidade, recomenda-se deixar sempre os '&' alinhados em todas
%as linhas e colunas

\chapter{Metodologia}

Aqui conterão os métodos e procedimentos adotados no desenvolvimento do
trabalho. Esta é uma das sessões mais importantes pois demonstra o poder
científico que foi utilizado para a pesquisa. Sem uma boa metodologia a
pesquisa pode perder a validade. O pesquisador deve utilizar métodos ou
técnicas aceitas pela comunidade científica na busca de provar suas hipóteses.

A metodologia escolhida deve ser aquela que mais se adéqua ao seu objeto de
estudo e à abordagem aplicada. Há dois métodos principais: 1) quantitativo, que
é o uso de instrumental estatístico, de dados numéricos; e 2) qualitativo, que
se caracteriza pela qualificação dos dados coletados, durante a análise do
problema.

\section{Uma seção}
Texto.

\section{Uma outra seção}
Texto.

\chapter{Apresentação e Análise dos Resultados}

Toda pesquisa deve apresentar uma análise sobre a investigação que foi
realizada através da metodologia que foi aplicada. Nesta sessão é interessante
inserir tabelas, gráficos, imagens que mostrem os resultados, análise de dados
coletados, etc.

É interessante que nessa sessão o autor compare os seus resultados com os
resultados de outros trabalhos existentes. Essa comparação aumenta a qualidade
do trabalho e demonstra a relevância do mesmo. 

\chapter{Conclusões e Trabalhos Futuros}

A conclusão deve conter os principais aspectos e contribuições de forma a
finalizar o trabalho apresentado. Deve-se apresentar o que era esperado do
trabalho através dos objetivos inseridos inicialmente e mostrar o que foi
conseguido.

Não deve-se inserir um novo assunto na conclusão. Aqui o autor apresentará as
próprias impressões sobre o trabalho efetuado.

É importante também que sejam identificadas limitações e problemas que surgiram
durante o desenvolvimento do trabalho e quais as consequências do mesmo.

Os trabalhos futuros devem conter oportunidades de expansão do trabalho
apresentado, bem como, novos projetos que puderam ser vislumbrados a partir do
desenvolvimento do trabalho


% Referencias

\begin{references}
  \bibliography{bib/references}
\end{references}

% Apendice

\theappendix
\chapter{Mapping Study's Instruments}
\label{ap:mapping-study}

\begin{table}[!htp]
	\centering
	\caption{List of conferences on which the searches were performed.}
	\label{tbl:conferences_list}
	\rowcolors{2}{lightgray!30}{white}
	\resizebox{\columnwidth}{!}{
	\begin{tabular}{ll}
	\toprule
	\textbf{Acronym} & \textbf{Conference} \\
	\toprule
	APSEC & Asia Pacific Software Engineering Conference \\
	ASE   & IEEE/ACM International Conference on Automated Software Engineering \\
	CSMR  & European Conference on Software Maintenance and Reengineering \\
	ESEC  & European Software Engineering Conference \\
	ESEM  & International Symposium on Empirical Software Management and Measurement \\
	ICSE  & International Conference on Software Engineering \\
	ICSM  & International Conference on Software Maintenance \\
	ICST & International Conference on Software Testing \\
	InfoVis & IEEE Information Visualization Conference \\
	KDD   & ACM SIGKDD International Conference on Knowledge Discovery and Data Mining \\
	MSR   & Working Conference on Mining Software Repositories \\
	OOPSLA & Object-Oriented Programming, Systems, Languages and Applications \\
	QSIC  & International Conference On Quality Software \\
	SAC & ACM Symposium on Applied Computing \\
	SEAA & EUROMICRO Conference on Software Engineering and Advanced Applications\\
	SEDE & 19th International Conference on Software Engineering and Data Engineering \\
	SEKE  & International Conference on Software Engineering and Knowledge Engineering \\
	\bottomrule
	\end{tabular}
	}
\end{table}

\begin{table}[htp]
	\caption{List of journals in which the searches were performed.}
	\label{tbl:journals_list}
	\centering
	\rowcolors{2}{lightgray!30}{white}
	\begin{tabular}{l}
	\toprule
	\textbf{Journal title} \\
	\toprule
	ACM Transactions on Software Engineering and Methodology \\
	Automated Software Engineering \\
	Elsevier Information and Software Technology \\
	Elsevier Journal of Systems and Software \\
	Empirical Software Engineering \\
	IEEE Software \\
	IEEE Computer \\
	IEEE Transactions on Software Engineering \\
	International Journal of Software Engineering and Knowledge Engineering \\
	Journal of Software: Evolution and Process \\
	Software Quality Journal \\
	Journal of Software \\
	Software Practice and Experience Journal \\
	\bottomrule
	\end{tabular}
\end{table}

\begin{table}[h]
\centering
\footnotesize
 \rowcolors{2}{lightgray!30}{white}
\caption{Search string per Search Engine.}
\label{tbl:stringengine}
\begin{tabular}{p{.15\textwidth}p{.8\textwidth}}
\toprule
\textbf{Search Engine} & \textbf{Search String}\\
\toprule
   	 	Google Scholar &  bug report OR track OR triage ``change
   	 	request'' issue track OR request OR software OR ``modification request'' OR
   	 	``defect track'' OR ``software issue''  repositories maintenance evolution\\

   	 	ACM Portal & Abstract: "bug report" or Abstract:"change request"
   	 	or Abstract:"bug track" or Abstract:"issue track" or  Abstract:"defect
   	 	track" or Abstract:"bug triage" or Abstract: "software issue" or Abstract: "issue request"
   	 	or Abstract: "modification request") and  (Abstract:software or
   	 	Abstract:maintenance or Abstract:repositories or Abstract:repository \\

   	 	IEEExplorer (1) & (((((((((("Abstract": "bug report") OR
   	 	"Abstract":"change request") OR "Abstract":"bug track") OR "Abstract":"software issue") OR "Abstract":"issue request") OR
        "Abstract":"modification request") OR "Abstract":"issue track") OR
	    "Abstract":"defect track") OR "Abstract":"bug triage") AND
	    "Abstract":software)\\

         IEEExplorer (2) & (((((((((("Abstract": "bug report") OR
         "Abstract":"change request") OR "Abstract":"bug track") OR "Abstract":"software issue") OR
         "Abstract":"issue request") OR "Abstract":"modification request") OR
         "Abstract":"issue track") OR "Abstract":"defect track") OR
         "Abstract":"bug triage") AND "Abstract":maintenance)\\

         IEEExplorer (3) & (((((((((("Abstract": "bug report") OR
         "Abstract":"change request") OR "Abstract":"bug track") OR "Abstract":"software issue") OR
         "Abstract":"issue request") OR "Abstract":"modification request") OR
         "Abstract":"issue track") OR "Abstract":"defect track") OR
         "Abstract":"bug triage") AND "Abstract":repositories)\\

         IEEExplorer & (((((((((("Abstract": "bug report") OR
         "Abstract":"change request") OR "Abstract":"bug track") OR "Abstract":"software issue") OR
         "Abstract":"issue request") OR "Abstract":"modification request") OR
         "Abstract":"issue track") OR "Abstract":"defect track") OR
         "Abstract":"bug triage") AND "Abstract": repository)\\

         Citeseer Library & (abstract: "bug report" OR abstract:"change request" OR abstract:"bug track" OR abstract:"issue track" OR
	     abstract:"defect track" OR abstract:"bug triage" OR abstract: "software
	     issue" OR abstract: "issue request" OR abstract: "modification request")
	     AND (abstract:software OR abstract:maintenance OR abstract:repositories OR
	     abstract:repository)\\

	     Elsevier & ("bug report" OR "change
	     request" OR "bug track" OR "issue track" OR "defect track" OR "bug triage" OR "software issue" OR  "issue request" OR
	    "modification request") AND (software OR maintenance OR repositories OR
	    repository)\\

	    Scirus & ("bug report" OR "change request" OR "bug track" OR "issue track" OR  "defect track" OR "bug triage" OR
        "software issue" OR  "issue request" OR "modification request") AND
	    (software maintenance OR repositories OR repository) ANDNOT (medical OR
	    aerospace)\\

	    ScienceDirect & ("bug report" OR "change request" OR "bug track"
	     OR "issue track" OR "defect track" OR "bug triage" OR "issue request" OR
	     "modification request") AND LIMIT-TO(topics, "soft ware")\\

	     Scopus & ("bug report" OR "change request" OR "bug track" OR
	     "issue track" OR  "defect track" OR "bug triage" OR "software issue" OR
	     "issue request" OR "modification request") AND (software maintenance OR
	     repositories OR repository)\\

	     Wiley & ("bug report" OR "change request"
	     OR "bug track" OR "issue track" OR  "defect track" OR "bug triage" OR
         "software issue" OR  "issue request" OR "modification request") AND
	     (software maintenance OR repositories OR repository)\\

	     ISI Web\newline of Knowledge & ("bug report" OR "change request" OR "bug
	     track" OR "issue track" OR  "defect track" OR "bug triage" OR "software issue" OR  "issue request" OR "modification request") AND
	    (software maintenance OR repositories OR repository) ANDNOT (medical OR
	    aerospace)\\

	    SpringerLink & ("bug report" OR "change request" OR "bug track" OR "issue track" OR  "defect track" OR "bug triage" OR
        "software issue" OR  "issue request" OR "modification request") AND
	    (software maintenance OR repositories OR repository) ANDNOT (medical OR
	    aerospace)\\
	\bottomrule
\end{tabular}
\end{table}

\end{document}
